\documentclass[11pt,letterpaper]{article}
\usepackage[T1]{fontenc}
\usepackage[utf8]{inputenc}
\usepackage{lmodern,geometry,microtype,xcolor,enumitem,hyperref,fancyhdr}
\geometry{margin=0.78in}
\definecolor{ink}{HTML}{17324D}
\definecolor{accent}{HTML}{126A78}
\hypersetup{colorlinks=true,linkcolor=accent,urlcolor=accent,breaklinks=true}
\setlength{\parindent}{0pt}
\setlength{\emergencystretch}{3em}
\setlist[itemize]{leftmargin=*,itemsep=3pt,topsep=3pt}
\pagestyle{fancy}\fancyhf{}\fancyhead[L]{\small Architecture Documentation Suite}\fancyhead[R]{\small v1.0}\fancyfoot[R]{\thepage}\setlength{\headheight}{15pt}
\begin{document}
\thispagestyle{empty}
{\color{accent}\large TECHNICAL REFERENCE \hfill NOVEMBER 2026}\par\vspace{1.1cm}
{\color{ink}\LARGE\bfseries\raggedright Architecture Documentation Curriculum}\par\vspace{0.35cm}
{\Large Eight modules, progressive labs and capstone assessment}\par\vspace{0.55cm}
\rule{\textwidth}{0.7pt}\par\smallskip
{\small Architecture Documentation Suite | Release 1.0}\par\vspace{0.7cm}
\tableofcontents
\newpage
\section*{1. Program design}\addcontentsline{toc}{section}{1. Program design}
Audience: software engineers and architects with repository, Git and basic UML familiarity. Duration: eight modules of six hours each (48 hours total), a proposed cadence. Each module allocates one hour concept study, one hour guided reading, three hours lab and one hour verification/review. Prerequisites: Git, Python basics, command line, ability to read service configuration. Tools: text editor, LaTeX, PlantUML, Python 3 and optionally Graphviz. Use an approved local repository; do not execute untrusted code. The source reading sequence follows V\&B chapters rather than reproducing book exercises.\par\smallskip
\section*{2. Module 1: Architecture, stakeholders and scope}\addcontentsline{toc}{section}{2. Module 1: Architecture, stakeholders and scope}
Reading: Prologue P.1--P.5 and chapter 9. Objectives: distinguish structure, view, style and description; identify at least five concerns; choose and justify viewpoints. Lab: inspect README and manifests of a small service, freeze baseline, create stakeholder/concern register, draft context model. Acceptance: every selected view maps to a concern, scope includes exclusions, baseline recorded. Deliverable: concern matrix and context packet.\par\smallskip
\section*{3. Module 2: Module structures}\addcontentsline{toc}{section}{3. Module 2: Module structures}
Reading: chapters 1--2 and Appendix A.2. Objectives: distinguish decomposition, uses and generalization; model layers and data responsibilities. Lab: trace imports and interfaces, create package diagram and dependency matrix, identify cycle or prove none within inspected scope. Acceptance: every observed relation has evidence, constraints explicitly distinguish allowed from observed dependencies. Deliverable: module view and catalog.\par\smallskip
\section*{4. Module 3: Runtime components and connectors}\addcontentsline{toc}{section}{4. Module 3: Runtime components and connectors}
Reading: chapters 3--4 and Appendix A.3. Objectives: distinguish static from runtime relations; characterize at least three connectors. Lab: trace a request and asynchronous or failure path; produce C4-style container and runtime component diagrams with contracts. Acceptance: all arrows labeled and every protocol marked known or unknown. Deliverable: C\&C packet and interface inventory.\par\smallskip
\section*{5. Module 4: Allocation and deployment}\addcontentsline{toc}{section}{5. Module 4: Allocation and deployment}
Reading: chapter 5 and Appendix A.4. Objectives: map artifacts to execution nodes and distinguish intended versus actual topology. Lab: inspect IaC or Compose, document deployment topology, configuration variants and ownership. Acceptance: every instance maps to a logical runtime unit or explicit exception; no live deployment claim without evidence. Deliverable: deployment and allocation mapping.\par\smallskip
\section*{6. Module 5: Interfaces, behavior and decisions}\addcontentsline{toc}{section}{6. Module 5: Interfaces, behavior and decisions}
Reading: chapters 6--8. Objectives: specify interfaces beyond signatures; represent happy path and error path; capture variation and rationale. Lab: create sequence and state diagrams, interface contract and two ADRs (observed versus proposed separated). Acceptance: failure, timeout, recovery and state termination addressed where relevant; decisions include alternatives and consequences. Deliverable: behavior/contract packets and ADRs.\par\smallskip
\section*{7. Module 6: Integrated framework and documentation package}\addcontentsline{toc}{section}{7. Module 6: Integrated framework and documentation package}
Reading: chapters 9--10, Epilogue E.1; official C4 and ISO public pages. Objectives: explain roles of V\&B, C4, UML, PlantUML and ISO; build cross-view correspondence. Lab: fill SAD template, construct view register, map requirements/concerns, render diagrams. Acceptance: no dangling IDs, all selected views include supporting documentation, standard claims labeled alignment only. Deliverable: integrated AD.\par\smallskip
\section*{8. Module 7: Agent harness and AppSec specialization}\addcontentsline{toc}{section}{8. Module 7: Agent harness and AppSec specialization}
Reading: chapters 3, 6--8 and suite document 07. Objectives: separate model proposal from authority; identify trust crossings, approval and memory threats; design denial and recovery tests. Lab: adapt supplied agent harness diagrams to a proposed system, model policy denial, scoped execution and checkpoint replay; create a security packet and threat register. Acceptance: every privileged crossing has enforcement or gap; no claim of implemented control without evidence. Deliverable: agent architecture description and test plan.\par\smallskip
\section*{9. Module 8: Review, release and capstone}\addcontentsline{toc}{section}{9. Module 8: Review, release and capstone}
Reading: chapter 11 and P.5. Objectives: evaluate documentation fitness, run gates and prepare change-controlled release. Lab: peer-review capstone with developer, operator and security roles; execute structural validator, render documents/diagrams, reconcile findings. Acceptance: all mandatory applicable local gates pass, no unresolved blocker, review records and unknowns included. Deliverable: release bundle, review report and remediation log.\par\smallskip
\section*{10. Assessment rubric}\addcontentsline{toc}{section}{10. Assessment rubric}
Use 100 points: scope/concerns 15; module/runtime/allocation modeling 25; interfaces/behavior/rationale 15; evidence and traceability 15; security and quality scenarios 15; review/build/readability 15. Grade each category by proportion of demonstrated requirements; incomplete or unsupported claims earn no credit. Passing threshold: 80/100 and no critical provenance or authorization-boundary defect. This is a proposed teaching rubric, not an accredited standard.\par\smallskip
\section*{11. Capstone progression and teaching notes}\addcontentsline{toc}{section}{11. Capstone progression and teaching notes}
Use the same repository across modules to avoid fragmented artifacts. At each review, ask students to locate a real claim, explain its evidence, distinguish static and runtime meanings and describe impact of a change. Provide feedback before expanding diagrams. Advanced extension: use an interoperable skill to generate a first-pass AD, then independently validate its claims and compare human versus agent output. Track false claims and missed concerns rather than measuring only document length. Source: V\&B chapter sequence and Appendix A; labs, schedule and rubric are original instructional design.\par\smallskip
\section*{Bibliographic foundation}
Paul Clements et al., \emph{Documenting Software Architectures: Views and Beyond}, 2nd ed., Addison-Wesley/Pearson, copyright 2011, ISBN 978-0-321-55268-6. Section and chapter references refer to the supplied edition. Additional framework sources are identified in the applicable documents.\par
\end{document}
